%%=============================================================================
%% INTRODUCCION
%% Titulo sin numerar, Arial 20 pt negrilla, 100 pt antes / 24 pt despues.
%%=============================================================================
\tituloPreliminar{Introducción}

Durante la última década, el Internet de las Cosas (Internet of Things, IoT) ha experimentado un crecimiento acelerado, transformando la forma en que los sistemas físicos interactúan con infraestructuras digitales mediante sensores, actuadores y dispositivos embebidos interconectados. Actualmente, el IoT constituye uno de los principales pilares de la denominada Industria 4.0 y soporta aplicaciones críticas en sectores como la salud, el transporte inteligente, la manufactura avanzada, la agricultura de precisión, las ciudades inteligentes y la infraestructura energética. La incorporación masiva de dispositivos conectados ha permitido incrementar los niveles de automatización, monitoreo en tiempo real y toma de decisiones basada en datos, favoreciendo procesos más eficientes y una mayor capacidad de respuesta frente a eventos dinámicos. Sin embargo, esta misma expansión ha incrementado significativamente la superficie de ataque disponible para actores maliciosos, convirtiendo la seguridad informática en uno de los principales desafíos tecnológicos de la actualidad.

Las proyecciones internacionales estiman que durante los próximos años existirán decenas de miles de millones de dispositivos IoT desplegados alrededor del mundo, muchos de ellos desempeñando funciones críticas relacionadas con infraestructura esencial, servicios médicos, redes eléctricas, transporte público e instalaciones industriales \cite{Sinha2025StateGlobally}. A diferencia de los sistemas informáticos tradicionales, estos dispositivos suelen operar durante largos periodos de tiempo sin supervisión directa, poseen capacidades limitadas de procesamiento, memoria y almacenamiento, y frecuentemente dependen de baterías o mecanismos de alimentación con restricciones energéticas severas. Estas características hacen que la implementación de mecanismos criptográficos convencionales resulte costosa o incluso inviable desde el punto de vista computacional, generando la necesidad de desarrollar soluciones específicamente diseñadas para ambientes con recursos restringidos \cite{McKay2017ReportCryptography,Thakor2021LightweightOpportunities}. Diversos estudios recientes coinciden en que la seguridad de los sistemas IoT debe abordarse mediante algoritmos capaces de ofrecer altos niveles de protección utilizando una cantidad mínima de memoria, ciclos de procesamiento y consumo energético \cite{Dhanda2020LightweightIoT,Seyhan2022ClassificationTaxonomy}.

Como respuesta a estas limitaciones surgió el área de investigación conocida como \textit{Lightweight Cryptography} (LWC), cuyo propósito consiste en diseñar mecanismos criptográficos capaces de preservar las propiedades fundamentales de confidencialidad, integridad, autenticidad y disponibilidad utilizando recursos considerablemente inferiores a los requeridos por algoritmos tradicionales \cite{McKay2017ReportCryptography}. La criptografía ligera no busca reemplazar completamente los estándares clásicos, sino proporcionar alternativas especialmente adaptadas a plataformas embebidas donde cada kilobyte de memoria, cada ciclo de reloj y cada miliwatt de energía representan recursos escasos. Durante la última década, esta área ha recibido un interés creciente tanto por parte de la comunidad científica como de organismos internacionales de estandarización, impulsando el desarrollo de nuevos algoritmos optimizados para microcontroladores, sensores inalámbricos y sistemas ciberfísicos \cite{Althobaiti2025LightweightStandardization,Radhakrishnan2024EfficiencyDevices}.

La importancia de esta línea de investigación llevó al National Institute of Standards and Technology (NIST) a iniciar en 2018 el proceso internacional de estandarización de algoritmos de criptografía ligera \cite{McKay2017ReportCryptography}. Después de evaluar decenas de propuestas provenientes de universidades, centros de investigación e industria, el proceso concluyó con la selección de ASCON como nuevo estándar para autenticación y cifrado autenticado en dispositivos con restricciones computacionales \cite{NationalInstituteofStandards2023AnnouncingSelection,Turan2025Ascon-BasedFunctions}. La elección de ASCON respondió a su capacidad para ofrecer elevados niveles de seguridad junto con un excelente desempeño tanto en implementaciones software como hardware, convirtiéndose actualmente en el principal referente para el desarrollo de aplicaciones IoT seguras \cite{Dobraunig2021AsconHashing}. No obstante, aun cuando ASCON proporciona mecanismos robustos para proteger la información transmitida, la seguridad global del sistema continúa dependiendo de la calidad del material criptográfico utilizado durante la generación de claves, semillas, vectores de inicialización y números aleatorios \cite{Katz2020IntroductionCryptography}.

La generación de números aleatorios constituye uno de los componentes más críticos dentro de cualquier arquitectura criptográfica moderna. Diversos ataques históricos han demostrado que incluso algoritmos considerados matemáticamente seguros pueden verse completamente comprometidos cuando utilizan fuentes de entropía deficientes o generadores de números pseudoaleatorios con baja calidad estadística \cite{Katz2020IntroductionCryptography,Menezes1996HandbookCryptography}. En consecuencia, la fortaleza de protocolos criptográficos, sistemas de autenticación, firmas digitales y mecanismos de establecimiento de claves depende directamente de la capacidad del sistema para producir secuencias impredecibles, no reproducibles y estadísticamente uniformes. Por esta razón, organismos como el NIST han desarrollado estándares específicos para la evaluación de fuentes de entropía, incluyendo las recomendaciones SP 800-22 para pruebas estadísticas, SP 800-90A para generadores determinísticos y SP 800-90B para estimación de entropía en fuentes físicas y no determinísticas \cite{Rukhin2010AApplications,Barker2015RecommendationGenerators,Turan2018RecommendationGeneration}.

Tradicionalmente, la obtención de entropía en dispositivos IoT ha dependido de generadores de números aleatorios implementados mediante hardware (\textit{True Random Number Generators}, TRNG), los cuales aprovechan fenómenos físicos como ruido térmico, osciladores de anillo, variaciones eléctricas o procesos cuánticos para producir secuencias impredecibles \cite{Seyhan2022ClassificationTaxonomy}. Aunque estos mecanismos ofrecen elevados niveles de seguridad, presentan limitaciones importantes relacionadas con el costo de implementación, la disponibilidad del hardware, la portabilidad entre plataformas y la complejidad de validación. En muchos dispositivos de bajo costo, especialmente aquellos destinados a aplicaciones masivas del Internet de las Cosas, la incorporación de hardware especializado resulta económicamente inviable, motivando la búsqueda de alternativas software capaces de generar niveles adecuados de entropía utilizando únicamente los recursos disponibles en el sistema operativo y la microarquitectura del procesador.

En este contexto surge \textit{Quantum Permutation Pad Random Number Generator} (QPP-RNG), una propuesta reciente desarrollada por Quantropi Inc. basada en la obtención de entropía mediante múltiples fuentes de comportamiento interno del sistema computacional \cite{Kuang2025QPP-RNG:Randomness,Kuang2025IID-BasedJitter}. A diferencia de los enfoques tradicionales, QPP-RNG combina información proveniente del \textit{system jitter}, tiempos de ejecución, acceso a memoria, variaciones microarquitecturales, comportamiento del planificador del sistema operativo y otros eventos no determinísticos para construir una fuente de entropía completamente implementada en software. Los resultados publicados muestran que este mecanismo alcanza niveles de min-entropía cercanos al máximo teórico permitido por byte y supera satisfactoriamente las principales pruebas estadísticas internacionalmente aceptadas, incluyendo NIST SP 800-22, NIST SP 800-90B y ENT, posicionándose como una alternativa prometedora para aplicaciones criptográficas de alta seguridad sin requerir hardware adicional \cite{Vrana2025RawEvaluation,Walker2008ENT:Program}.

Paralelamente, el avance acelerado de la computación cuántica ha introducido nuevos desafíos para la seguridad de los sistemas criptográficos actuales. Algoritmos como el propuesto por Shor demuestran que los principales esquemas de criptografía asimétrica utilizados actualmente, entre ellos RSA y la criptografía de curva elíptica (ECC), podrán ser vulnerables una vez existan computadores cuánticos suficientemente grandes y tolerantes a fallos \cite{Shor1994AlgorithmsFactoring,Ivezic2017ShorsCryptography}. Aunque la criptografía ligera se concentra principalmente en algoritmos simétricos, la transición hacia ecosistemas resistentes a la computación cuántica requiere fortalecer todos los componentes de la arquitectura criptográfica, incluyendo la generación de material aleatorio de alta calidad. Dentro de este escenario, Quantum Permutation Pad (QPP) ha sido propuesto como un marco criptográfico capaz de aprovechar propiedades derivadas de espacios de permutaciones para construir mecanismos resistentes frente a ataques clásicos y cuánticos, ofreciendo además niveles significativamente superiores de entropía respecto a enfoques tradicionales \cite{Kuang2022QuantumCryptography,Kuang2020ShannonSpace}.

A pesar de los resultados obtenidos por QPP y QPP-RNG, actualmente existe una brecha importante entre las propuestas teóricas disponibles en la literatura y su utilización práctica dentro de dispositivos IoT críticos. En particular, no se dispone de implementaciones abiertas, modulares y optimizadas que permitan integrar QPP-RNG como fuente de entropía dentro de arquitecturas modernas de criptografía ligera. Asimismo, existe una limitada evidencia experimental respecto al comportamiento de este mecanismo sobre plataformas embebidas, especialmente en términos de consumo de memoria, utilización del procesador, rendimiento, calidad estadística y facilidad de integración con algoritmos ligeros ampliamente utilizados como ASCON.

En este contexto, la presente investigación propone el diseño, implementación y evaluación de una librería criptográfica ligera desarrollada en el lenguaje de programación Rust, incorporando una adaptación optimizada de QPP-RNG como fuente de entropía para dispositivos IoT críticos. Rust fue seleccionado debido a sus mecanismos de seguridad de memoria, ausencia de \textit{data races}, modelo de propiedad (\textit{ownership}) y creciente adopción dentro del desarrollo de software seguro para infraestructura crítica. La investigación contempla el análisis integral de la calidad estadística de los números aleatorios generados, el estudio del consumo de recursos computacionales, la evaluación del desempeño frente a soluciones de referencia y la validación de la viabilidad práctica de incorporar QPP-RNG dentro de arquitecturas de criptografía ligera orientadas a sistemas embebidos.

Finalmente, esta investigación busca contribuir simultáneamente en tres dimensiones. Desde la perspectiva científica, amplía el conocimiento existente sobre la utilización de nuevas fuentes de entropía software para criptografía ligera y seguridad post-cuántica. Desde el punto de vista tecnológico, desarrolla una librería modular reutilizable que podrá servir como base para futuras investigaciones y aplicaciones industriales. Finalmente, desde una perspectiva práctica, proporciona evidencia experimental acerca de la viabilidad de incorporar mecanismos avanzados de generación de entropía dentro de dispositivos IoT con restricciones severas de memoria, procesamiento y consumo energético, contribuyendo al fortalecimiento de la seguridad de la próxima generación de sistemas inteligentes conectados.